\scanpage{75}
% archive foliation: 74; typescript printed page number follows
\begin{center}
- 19 -
\end{center}

On notera que tout système de $q$ séries entières $u_i$ sans terme constant
$(1\leq i\leq q)$ sommables dans un voisinage de l'origine, est substituable
dans chacune des séries entières
$\sum_{(n_j)}h_{i;n_1\ldots n_q}y_1^{n_1}\ldots y_q^{n_q}$
(cor. du th.2) ; si le problème a une solution, les séries entières
(par rapport aux $x_i$) des deux membres de (5) sont donc sommables au voisinage
de l'origine, et par suite doivent avoir leurs coefficients \uline{identiques}.

Supposons donc le problème résolu, et désignons par
$c_i(n_1,\ldots,\unsure{q|n_q};r_1,\ldots,r_p)$ le coefficient de
$x_1^{r_1}\ldots x_p^{r_p}$ dans la série $h_{i;n_1\ldots n_q}$,
par $c_i(r_1,\ldots,r_p)$ le coefficient du même terme dans $h_{i0}$
par $d_i(m_1,\ldots,m_p)$ le coefficient de $x_1^{m_1}\ldots x_p^{m_p}$
dans $u_i$. L'identification des séries des deux membres de (5) donne les conditions
\[
(6)\quad
\left\{\begin{aligned}
d_i(0,0,\ldots,0)&=0\\
d_i(m_1,\ldots,m_p)&=c_i(m_1,\ldots,m_p)+\\
&+\sum_{(n_j)(r_i)}
\left(c_i(n_1,\ldots,n_q;r_1,\ldots,r_p)
\prod_{j=1}^q\prod_{k=1}^{n_j}d_j(t_{kj1},t_{kj2},\ldots,t_{kjp})\right)
\end{aligned}\right.
\]
la sommation du second membre étant étendue à tous les systèmes d'entiers
$\geq0$, $n_j,r_i$ et $t_{kji}$ satisfaisant aux conditions
$\sum_{i=1}^p t_{kji}\geq1$, $\sum_{j=1}^q n_j\geq2$ et
\[
(7)\qquad
m_i=r_i+\sum_{j=1}^q\sum_{k=1}^{n_j}t_{kji}
\qquad(1\leq i\leq p).
\]
On tire de ces conditions que
\[
\sum_{i=1}^p m_i=\sum_{i=1}^p r_i+
\sum_{j=1}^q\sum_{k=1}^{n_j}\sum_{i=1}^p t_{kji}
\]
Si on tient compte de ce que $\sum_{i=1}^p t_{kji}\geq1$
et $\sum_{j=1}^q n_j\geq2$, on voit que pour tout système
$(t_{kji})_{1\leq i\leq p}$ qui figure au second membre de (7),
on a $\sum_{i=1}^p t_{kji}<\sum_{i=1}^p m_i$, en d'autres termes,
les relations (6) déterminent les coefficients des termes de degré
\uline{total} $s$ dans $u_i$ en fonction de ceux des termes de degré total
$<s$ dans les $u_j$ $(1\leq j\leq q)$ ; par récurrence sur $s$,
on voit donc que les équations (6) déterminent les coefficients
